% ROZDZIAŁ 22 — ZBIEŻNOŚĆ MIAR I PROCHOROW
\chapter[Zbieżność miar i twierdzenie Prochorowa]{Zbieżność miar, twierdzenie Portmanteau i zwartość Prochorowa}
\label{ch:weak-measures-22}

\section{Trzy topologie, których nie należy utożsamiać}
Niech $(X,d)$ będzie przestrzenią metryczną,
$\mathcal B(X)$ jej $\sigma$-algebrą borelowską,
a $\mathcal P(X)$ zbiorem miar probabilistycznych na $X$.
Przestrzeń $C_b(X)$ składa się z ograniczonych rzeczywistych
funkcji ciągłych. Jeśli $X$ jest dodatkowo lokalnie zwarta,
$C_c(X)$ oznacza przestrzeń zdefiniowaną w
rozdziale~\ref{ch:radon-measures-20}, a $C_0(X)$
przestrzeń wprowadzoną w rozdziale~\ref{ch:riesz-markov-21}.

\begin{definicja}[Zbieżność wąska]
\label{def:narrow-measures-22}
Ciąg $(\mu_n)\subseteq\mathcal P(X)$ jest
\emph{wąsko zbieżny} do $\mu\in\mathcal P(X)$,
co zapisujemy $\mu_n\Rightarrow\mu$, jeśli
\begin{equation}\label{eq:narrow-22}
 \int_X f(x)\,\mathrm d\mu_n(x)
 \longrightarrow\int_X f(x)\,\mathrm d\mu(x)
 \quad\text{dla każdej }f\in C_b(X).
\end{equation}
\end{definicja}

\begin{definicja}[Zbieżność słaba-* i mglista]
\label{def:vague-weakstar-22}
Niech $X$ będzie lokalnie zwartą ośrodkową przestrzenią
metryczną, a $\mu_n,\mu$ miarami Radona.
Nazywamy zbieżność $\mu_n\to\mu$ \emph{mglistą},
gdy $\int_Xf\,\mathrm d\mu_n\to\int_Xf\,\mathrm d\mu$
dla każdego $f\in C_c(X)$.
Jeżeli miary są skończone, zbieżność
$\sigma(C_0(X)^*,C_0(X))$ (\emph{słaba-*})
oznacza zbieżność całek dla wszystkich $f\in C_0(X)$.
W przypadku miar probabilistycznych zbieżność wąska
implikuje obie te zbieżności. Na przestrzeni niezwartej
implikacje odwrotne nie muszą zachodzić.
\end{definicja}

\begin{przyklad}[Ucieczka masy]
\label{ex:mass-escape-22}
Niech $X=\mathbb R$, $\mu_n=\delta_n$.
Dla każdego $f\in C_0(\mathbb R)$ zachodzi $f(n)\to0$,
zatem $\mu_n\to0$ słabo-* w dualnej przestrzeni
$C_0(\mathbb R)^*$. Także $\mu_n\to0$ mgliście.
Jednak $\mu_n$ nie zbiega wąsko do żadnej miary
probabilistycznej: stała funkcja $1\in C_b(\mathbb R)$
ma całkę równą $1$ względem każdej $\mu_n$.
W szczególności określenie „zbieżność słaba miar” bez
wskazania klasy funkcji testowych bywa dwuznaczne.
\end{przyklad}

\section{Twierdzenie Portmanteau}
\begin{definicja}[Zbiór ciągłości miary]
\label{def:continuity-set-22}
Niech $\mu\in\mathcal P(X)$. Zbiór borelowski
$A\subseteq X$ nazywamy \emph{zbiorem ciągłości miary}
$\mu$, jeżeli $\mu(\partial A)=0$, gdzie
$\partial A=\overline A\setminus A^\circ$ jest jego brzegiem.
\end{definicja}

\begin{twierdzenie}[Portmanteau]
\label{thm:portmanteau-22}
Niech $(X,d)$ będzie dowolną przestrzenią metryczną,
a $\mu_n,\mu\in\mathcal P(X)$. Następujące warunki
są równoważne:
\begin{enumerate}[label=(\roman*)]
\item $\mu_n\Rightarrow\mu$;
\item dla każdej funkcji ograniczonej, dolnie półciągłej
$f:X\to\mathbb R$ zachodzi
$\int_Xf\,\mathrm d\mu\le\liminf_n\int_Xf\,\mathrm d\mu_n$;
\item dla każdego otwartego $G\subseteq X$ zachodzi
$\mu(G)\le\liminf_n\mu_n(G)$;
\item dla każdego domkniętego $F\subseteq X$ zachodzi
$\limsup_n\mu_n(F)\le\mu(F)$;
\item dla każdego zbioru ciągłości $A$ miary $\mu$
zachodzi $\mu_n(A)\to\mu(A)$.
\end{enumerate}
W warunku (ii) można także dopuścić dolnie półciągłe
funkcje ograniczone z dołu, przy całkach w $(-\infty,+\infty]$.
\end{twierdzenie}
\begin{proof}
\emph{(i)$\Rightarrow$(iii).}
Dla otwartego $G$ określamy
$f_k(x)=\min\{1,kd(x,X\setminus G)\}$, gdy
$G\neq X$, oraz $f_k\equiv1$, gdy $G=X$.
Funkcje $f_k$ są ciągłe, $0\le f_k\le\chi_G$,
a $f_k\uparrow\chi_G$. Zatem dla każdego $k$
\[
\liminf_n\mu_n(G)\ge\lim_n\int_X f_k\,\mathrm d\mu_n
=\int_X f_k\,\mathrm d\mu.
\]
Z twierdzenia Beppo Leviego, po przejściu $k\to\infty$,
otrzymujemy (iii).

\emph{(iii)$\Leftrightarrow$(iv).}
Dla $F$ domkniętego zbiór $X\setminus F$ jest otwarty,
a $\mu_n(X\setminus F)=1-\mu_n(F)$ i podobnie
dla $\mu$. Pozwala to zamienić dolną granicę na górną.

\emph{(iii)$\Rightarrow$(ii).}
Najpierw niech $f$ będzie dolnie półciągła i
$0\le f\le M$. Dla każdego $t\in[0,M]$ zbiór
$G_t=\{x:f(x)>t\}$ jest otwarty. Z warunku (iii)
$\mu(G_t)\le\liminf_n\mu_n(G_t)$.
Z reprezentacji warstwowej całki
$\int_X f\,\mathrm d\mu=\int_{[0,M]}\mu(G_t)
\,\mathrm d\lambda_1(t)$, twierdzenia Tonellego
i lematu Fatou otrzymujemy
\[
\int_X f\,\mathrm d\mu
\le\int_{[0,M]}\liminf_n\mu_n(G_t)\,\mathrm d\lambda_1(t)
\le\liminf_n\int_Xf\,\mathrm d\mu_n.
\]
Jeśli $f\ge a$ i nie jest ograniczona z góry,
stosujemy poprzedni argument do
$f_k=\min\{f-a,k\}$, a następnie Beppo Leviego;
odjęta stała $a$ ma ten sam całkowy wkład dla
wszystkich miar probabilistycznych.

\emph{(ii)$\Rightarrow$(i).}
Dla dowolnej $h\in C_b(X)$ zarówno $h$, jak i $-h$
są dolnie półciągłe i ograniczone z dołu. Stąd
$\int h\,\mathrm d\mu\le\liminf_n\int h\,\mathrm d\mu_n$
oraz $-\int h\,\mathrm d\mu\le
-\limsup_n\int h\,\mathrm d\mu_n$.
Obie nierówności wymuszają zbieżność całek.

\emph{(iii), (iv)$\Rightarrow$(v).}
Dla $A$ spełniającego $\mu(\partial A)=0$ mamy
$A^\circ\subseteq A\subseteq\overline A$ oraz
$\mu(A^\circ)=\mu(A)=\mu(\overline A)$.
Wobec tego
\[
\mu(A)\le\liminf_n\mu_n(A^\circ)
\le\liminf_n\mu_n(A)
\le\limsup_n\mu_n(A)
\le\limsup_n\mu_n(\overline A)\le\mu(A).
\]

\emph{(v)$\Rightarrow$(i).}
Niech $h\in C_b(X)$, $\|h\|_\infty\le M$.
Miara obrazowa $h_*\mu$ na $[-M,M]$ ma co najwyżej
przeliczalnie wiele atomów: dla każdego $k$ liczba
atomów o masie co najmniej $1/k$ jest skończona.
Zatem można wybrać podział
$-M=t_0<\cdots<t_N=M$ o dostatecznie małych oczkach,
którego punkty wewnętrzne są nieatomowe dla $h_*\mu$.
Aby uniknąć atomów w końcach, wolno wybrać
$t_0<-M$ i $t_N>M$ oraz wymagać, by również one
były nieatomowe. Dla przedziałów $(t_{j-1},t_j]$
zbiory $A_j=\{t_{j-1}<h\le t_j\}$ mają
$\partial A_j\subseteq\{h=t_{j-1}\}\cup\{h=t_j\}$,
więc są zbiorami ciągłości miary. Z (v) wynika
$\mu_n(A_j)\to\mu(A_j)$. Aproksymując $h$
na każdym $A_j$ przez stałą wybraną z przedziału,
z błędem co najwyżej długość przedziału, otrzymujemy
\[
 \limsup_n\left|\int_Xh\,\mathrm d\mu_n-
 \int_Xh\,\mathrm d\mu\right|
 \le2\max_j(t_j-t_{j-1}).
\]
Wybierając dowolnie drobne podziały kończymy dowód.
\end{proof}

\begin{przyklad}[Dlaczego przedziały nie zawsze są zbiorami ciągłości?]
Dla $\mu_n=\delta_{1/n}$ i $\mu=\delta_0$ na
$\mathbb R$ mamy $\mu_n\Rightarrow\mu$.
Tymczasem $\mu_n(\{0\})=0$ i $\mu(\{0\})=1$.
Brzeg zbioru $\{0\}$ ma dodatnią miarę $\mu$,
więc Portmanteau nie gwarantuje zbieżności
wartości miar tego zbioru.
\end{przyklad}

\begin{wniosek}[Kiedy zbieżność mglista jest słabą-*?]
\label{cor:vague-c0-bounded-22}
Niech $X$ będzie lokalnie zwartą ośrodkową przestrzenią
metryczną i niech $\mu_n,\mu$ będą skończonymi dodatnimi
miarami Radona. Jeśli $\sup_n\mu_n(X)<\infty$, to
zbieżność mglista $\mu_n\to\mu$ jest równoważna zbieżności
$\sigma(C_0(X)^*,C_0(X))$.
Bez ograniczenia mas implikacja od zbieżności mglistej
do słabej-* może zawieść.
\end{wniosek}
\begin{proof}
Zbieżność słaba-* daje mglistą, gdyż $C_c(X)\subset C_0(X)$.
Załóżmy zbieżność mglistą oraz $C:=\sup_n\mu_n(X)<\infty$.
Niech $f\in C_0(X)$ i $\varepsilon>0$. Z gęstości
$C_c(X)$ w $C_0(X)$ (lemat~\ref{lem:density-cc-c0-21})
wybieramy $g\in C_c(X)$ takie, że
$\|f-g\|_\infty<\varepsilon$. Wówczas
\[
 \left|\int_X f\,\mathrm d\mu_n-\int_X f\,\mathrm d\mu\right|
 \le\left|\int_Xg\,\mathrm d\mu_n-\int_Xg\,\mathrm d\mu\right|
 +(C+\mu(X))\varepsilon.
\]
Pierwszy składnik dąży do zera, a $\varepsilon$ jest dowolne.

Aby wykazać konieczność kontroli mas, weźmy
$X=\mathbb R$, $\mu_n=n\delta_n$ i $\mu=0$.
Dla każdego $g\in C_c(\mathbb R)$ zachodzi
$\int g\,\mathrm d\mu_n=0$ dla dostatecznie dużych $n$,
więc zachodzi zbieżność mglista. Funkcja
$f(x)=(1+|x|)^{-1/2}$ należy do $C_0(\mathbb R)$,
a $\int f\,\mathrm d\mu_n=n/\sqrt{1+n}\to+\infty$.
Zbieżności słabej-* nie ma.
\end{proof}

\section{Zbieżność względem odległości w wariacji całkowitej}
\begin{twierdzenie}[Zbieżność w wariacji implikuje wąską]
\label{thm:tv-implies-narrow-22}
Jeśli $\mu_n,\mu\in\mathcal P(X)$ i
$d_{\mathrm{TV}}(\mu_n,\mu)\to0$, to
$\mu_n\Rightarrow\mu$.
\end{twierdzenie}
\begin{proof}
Dla dowolnej $h\in C_b(X)$, korzystając z normy
wariacji miary ze znakiem, mamy
\[
\left|\int_Xh\,\mathrm d\mu_n-\int_Xh\,\mathrm d\mu\right|
\le \|h\|_\infty\|\mu_n-\mu\|_{\mathrm{TV}}
=2\|h\|_\infty d_{\mathrm{TV}}(\mu_n,\mu)\to0.
\]
\end{proof}

\begin{uwaga}[Brak implikacji odwrotnej]
Przykład $\delta_{1/n}\Rightarrow\delta_0$ pokazuje,
że wąska zbieżność nie implikuje zbieżności w wariacji:
$d_{\mathrm{TV}}(\delta_{1/n},\delta_0)=1$.
\end{uwaga}

\section{Ciasność rodzin miar}
\begin{definicja}[Jednostajna ciasność]
\label{def:uniform-tight-22}
Niech $\mathcal F\subseteq\mathcal P(X)$.
Rodzina $\mathcal F$ jest \emph{jednostajnie ciasna},
jeżeli dla każdego $\varepsilon>0$ istnieje
zwarty zbiór $K\subseteq X$ spełniający
$\inf_{\mu\in\mathcal F}\mu(K)>1-\varepsilon$.
\end{definicja}

\begin{lemat}[Przybliżanie lipschitzowskie i przedłużanie]
\label{lem:lipschitz-extend-22}
Niech $(Y,D)$ będzie przestrzenią metryczną, $K\subseteq Y$
niepustym zbiorem zwartym, a $h:K\to\mathbb R$ funkcją
ciągłą z $\|h\|_\infty\le M$. Dla każdego $\eta>0$
istnieje ograniczona funkcja Lipschitzowska
$g:Y\to\mathbb R$ spełniająca
$\|g\|_\infty\le M$ oraz
$\sup_{x\in K}|g(x)-h(x)|<\eta$.
\end{lemat}
\begin{proof}
Funkcja $h$ jest jednostajnie ciągła na $K$.
Dla $L>0$ definiujemy
$h_L(x)=\inf_{z\in K}(h(z)+L D(x,z))$, $x\in K$.
Z nierówności trójkąta wynika
$|h_L(x)-h_L(y)|\le L D(x,y)$; ponadto
$h_L\le h$ i $-M\le h_L\le M$.
Wybierzmy $\delta>0$ takie, że
$|h(z)-h(x)|<\eta/2$ dla $D(x,z)<\delta$.
Jeżeli $L>2M/\delta$, to dla $D(x,z)\ge\delta$
mamy $h(z)+LD(x,z)\ge M\ge h(x)$.
Dla $D(x,z)<\delta$ natomiast
$h(z)+LD(x,z)\ge h(x)-\eta/2$.
Zatem $0\le h-h_L\le\eta/2$ na $K$.

Funkcję $h_L$ przedłużamy na $Y$ wzorem
$G(y)=\inf_{z\in K}(h_L(z)+L D(y,z))$.
Ponownie $G$ jest $L$-lipschitzowska. Na $K$
mamy $G=h_L$: nierówność $G\le h_L$ wynika
z podstawienia $z=y$, a odwrotna z własności
Lipschitza funkcji $h_L$. Funkcja
$g(y)=\max\{-M,\min\{M,G(y)\}\}$ jest Lipschitzowska,
ma normę co najwyżej $M$ i na $K$ jest równa $h_L$.
Wobec tego $\sup_K|g-h|\le\eta/2<\eta$.
\end{proof}

\begin{lemat}[Zwartość prawdopodobieństw na przestrzeni zwartej]
\label{lem:compact-prob-22}
Jeśli $Y$ jest zwartą przestrzenią metryczną,
to każdy ciąg w $\mathcal P(Y)$ ma podciąg wąsko
zbieżny do pewnego elementu $\mathcal P(Y)$.
\end{lemat}
\begin{proof}
Przestrzeń $C(Y)$ jest separowalna w normie supremum.
Aby to udowodnić bez odwołania do twierdzenia
Stone'a--Weierstrassa, wybierzmy przeliczalny zbiór
gęsty $(y_j)$ i rozważmy przeliczalną rodzinę
funkcji postaci
$F(x)=\min_{1\le i\le N}\{q_i+L d(x,y_{j_i})\}$,
gdzie $N,j_i\in\mathbb N$, $q_i\in\mathbb Q$,
$L\in\mathbb Q_{>0}$.
Pokażemy, że jest ona gęsta. Dla dowolnej
$h\in C(Y)$ i $\varepsilon>0$ lemat
\ref{lem:lipschitz-extend-22} (z $K=Y$) daje
funkcję $L$-lipschitzowską $v$ z
$\|h-v\|_\infty<\varepsilon/3$.
Zwiększając $L$, możemy przyjąć go wymiernym.
Wybieramy skończoną $\delta$-siatkę
$y_{j_1},\ldots,y_{j_N}$ oraz wymierne $q_i$
bliskie $v(y_{j_i})$, tak by
$|q_i-v(y_{j_i})|<\varepsilon/6$ i
$2L\delta<\varepsilon/6$.
Z lipschitzowskości wynika
$F(x)\ge v(x)-\varepsilon/6$ oraz
$F(x)\le v(x)+\varepsilon/6+2L\delta$.
W konsekwencji $\|F-h\|_\infty<\varepsilon$.
Niech $(f_j)$ będzie taką gęstą
rodziną. Dla danego ciągu $(\mu_n)$ każdy ciąg
$\int_Y f_j\,\mathrm d\mu_n$ jest ograniczona.
Metoda diagonalna dostarcza podciągu $(\mu_{n_k})$,
dla którego wszystkie te całki mają granice.
Ponieważ
$|\int h\,\mathrm d\mu_{n_k}-\int f_j\,\mathrm d\mu_{n_k}|
\le\|h-f_j\|_\infty$, te granice rozszerzają się
jednoznacznie na ciągły dodatni funkcjonał $L$ na
$C(Y)$ o normie $1$ i $L(\chi_Y)=1$.
Na mocy twierdzenia~\ref{thm:compact-riesz-21}
istnieje prawdopodobieństwo $\mu$ reprezentujące $L$.
Definicja $L$ daje $\mu_{n_k}\Rightarrow\mu$.
\end{proof}

\begin{lemat}[Włożenie separowalnej przestrzeni metrycznej
w zwartą przestrzeń metryczną]
\label{lem:hilbert-cube-22}
Dla każdej separowalnej przestrzeni metrycznej $(X,d)$
istnieje zwarta przestrzeń metryczna $Y$ i homeomorfizm
$\iota:X\to\iota(X)\subseteq Y$.
\end{lemat}
\begin{proof}
Niech $(x_j)$ będzie przeliczalnym zbiorem gęstym.
Definiujemy
\[
\iota(x)=\left(\frac{d(x,x_j)}{1+d(x,x_j)}\right)_{j\ge1}
\in[0,1]^{\mathbb N}.
\]
Przestrzeń kostki Hilberta $[0,1]^{\mathbb N}$ z metryką
$D(a,b)=\sum_{j\ge1}2^{-j}|a_j-b_j|$ jest zwarta:
metoda diagonalna zapewnia zbieżny podciąg współrzędnych,
a oszacowanie ogona sumy daje zbieżność w $D$.
Każda współrzędna $\iota$ jest ciągła, więc
$\iota$ jest ciągła. Jeśli $\iota(x)=\iota(y)$,
to $d(x,x_j)=d(y,x_j)$ dla wszystkich $j$;
przez gęstość wynika $x=y$.
Aby wykazać ciągłość odwrotności na obrazie, załóżmy,
że $\iota(x_n)\to\iota(x)$, i ustalmy $\varepsilon>0$.
Wybierzmy $x_j$ z $d(x,x_j)<\varepsilon/3$.
Dla dużych $n$ współrzędna $j$ daje
$|d(x_n,x_j)-d(x,x_j)|<\varepsilon/3$,
zatem $d(x_n,x)\le d(x_n,x_j)+d(x_j,x)<\varepsilon$.
Wystarczy przyjąć $Y=\overline{\iota(X)}$ w kostce.
\end{proof}

\begin{twierdzenie}[Prochorow: kierunek ciasność $\Rightarrow$ zwartość]
\label{thm:prokhorov-forward-22}
Niech $X$ będzie separowalną przestrzenią metryczną,
a $\mathcal F\subseteq\mathcal P(X)$ rodziną jednostajnie
ciasną. Wtedy każdy ciąg w $\mathcal F$ ma podciąg
wąsko zbieżny do prawdopodobieństwa na $X$.
\end{twierdzenie}
\begin{proof}
Niech $(\mu_n)\subseteq\mathcal F$.
Z jednostajnej ciasności wybieramy rosnący ciąg
zwartych zbiorów $K_m\subseteq X$ taki, że
\begin{equation}\label{eq:tight-compacts-22}
 \inf_n\mu_n(K_m)\ge1-2^{-m},\qquad m\ge1.
\end{equation}
Stosujemy włożenie $\iota:X\to Y$ z lematu
\ref{lem:hilbert-cube-22}. Miary
$\nu_n=\iota_*\mu_n$ są probabilistyczne na
zwartej metrycznej $Y$. Z lematu
\ref{lem:compact-prob-22} wybieramy podciąg,
ponownie oznaczany $(\nu_n)$, wąsko zbieżny
do $\nu\in\mathcal P(Y)$.
Dla każdego $m$ zbiór $\iota(K_m)$ jest zwarty,
a zatem domknięty w $Y$. Z Portmanteau
\[
\nu(\iota(K_m))\ge\limsup_n\nu_n(\iota(K_m))
\ge1-2^{-m}.
\]
Wobec rosnącości $K_m$ zbiór
$C=\bigcup_m\iota(K_m)$ jest borelowski w $Y$,
a $\nu(C)=1$.

Definiujemy teraz prawdopodobieństwo $\mu$ na $X$.
Dla $A\in\mathcal B(X)$ kładziemy
\begin{equation}\label{eq:pullback-limit-22}
\mu(A)=\nu\bigl(\iota(A\cap\bigcup_mK_m)\bigr).
\end{equation}
Ta definicja jest poprawna: $A\cap K_m$ jest
borelowski w zwartej przestrzeni $K_m$,
a $\iota|_{K_m}$ jest homeomorfizmem na zwarty
podzbiór $Y$. Zatem $\iota(A\cap K_m)$ jest
borelowski w $Y$, podobnie jak ich przeliczalna suma.
Iniektywność $\iota$ i addytywność $\nu$ dowodzą,
że $\mu$ jest miarą; ponadto $\mu(X)=\nu(C)=1$
i $\mu(K_m)=\nu(\iota(K_m))$.

Niech $f\in C_b(X)$, $M=\|f\|_\infty$.
Dla ustalonego $m$ funkcja $f\circ\iota^{-1}$
jest ciągła na zwartym $\iota(K_m)$.
Lemat~\ref{lem:lipschitz-extend-22}, zastosowany
z dowolnym $\eta>0$, daje funkcję
$g_{m,\eta}\in C_b(Y)$, $\|g_{m,\eta}\|_\infty\le M$,
która na $\iota(K_m)$ różni się od
$f\circ\iota^{-1}$ o mniej niż $\eta$.
Stąd, z \eqref{eq:tight-compacts-22},
\[
\left|\int_X f\,\mathrm d\mu_n-
\int_Y g_{m,\eta}\,\mathrm d\nu_n\right|\le2M2^{-m}+\eta,
\]
a przez analogiczny rachunek
$|\int_Xf\,\mathrm d\mu-\int_Yg_{m,\eta}\,\mathrm d\nu|
\le2M2^{-m}+\eta$.
Dla ustalonych $m,\eta$ granica całek z $g_{m,\eta}$ wynika
ze zbieżności wąskiej na $Y$; dlatego
\[
\limsup_n\left|\int_X f\,\mathrm d\mu_n-
\int_X f\,\mathrm d\mu\right|\le4M2^{-m}+2\eta.
\]
Przechodząc najpierw $\eta\to0$, a następnie
$m\to\infty$, kończymy dowód.
\end{proof}

\begin{lemat}[Jednostajna ciasność ciągu wąsko zbieżnego]
\label{lem:weak-sequence-tight-22}
Jeśli $X$ jest polską przestrzenią metryczną,
$\mu_n\Rightarrow\mu$ i $\mu_n,\mu\in\mathcal P(X)$,
to rodzina $\{\mu,\mu_1,\mu_2,\ldots\}$ jest
jednostajnie ciasna.
\end{lemat}
\begin{proof}
Niech $d$ będzie zupełną, zgodną metryką, a $(x_j)$
przeliczalnym zbiorem gęstym. Ustalmy $\varepsilon>0$.
Dla każdego $m$ przeliczalna rodzina kul otwartych
$B(x_j,2^{-m})$ pokrywa $X$. Istnieje skończona
suma tych kul $O_m$, dla której
$\mu(O_m)>1-\varepsilon2^{-m-2}$.
Z otwartej wersji Portmanteau wynika, że
$\mu_n(O_m)>1-\varepsilon2^{-m-1}$ dla wszystkich
$n$ od pewnego $N_m$.
Powiększając skończoną rodzinę kul w $O_m$,
możemy zapewnić tę samą nierówność dla
$\mu_1,\ldots,\mu_{N_m-1}$; pozostaje prawdziwa
także dla wszystkich późniejszych miar.
Niech $F_m$ będzie sumą domkniętych kul
$\overline B(x_j,2^{-m})$ o środkach użytych w $O_m$.
Kładziemy $K=\bigcap_mF_m$. Zbiór $K$ jest domknięty,
a dla każdego $m$ mieści się w skończonej sumie kul
o promieniu $2^{-m}$, jest zatem całkowicie ograniczony.
Zupełność $X$ daje zwartość $K$. Ponadto
$\mu_n(X\setminus K)\le\sum_{m\ge1}\mu_n(X\setminus F_m)
<\sum_{m\ge1}\varepsilon2^{-m-1}<\varepsilon$
dla wszystkich $n$. Podobne oszacowanie zachodzi dla $\mu$.
\end{proof}

\begin{twierdzenie}[Prochorow na przestrzeni polskiej]
\label{thm:prokhorov-22}
Niech $X$ będzie przestrzenią polską, a
$\mathcal F\subseteq\mathcal P(X)$.
Następujące warunki są równoważne:
\begin{enumerate}[label=(\roman*)]
\item $\mathcal F$ jest jednostajnie ciasna;
\item każdy ciąg w $\mathcal F$ ma podciąg wąsko
zbieżny do pewnego prawdopodobieństwa na $X$.
\end{enumerate}
Warunek (ii) oznacza względną zwartość
ciągową dla topologii zbieżności wąskiej.
\end{twierdzenie}
\begin{proof}
Implikację (i)$\Rightarrow$(ii) udowodniliśmy w twierdzeniu
\ref{thm:prokhorov-forward-22}. Dla odwrotnej
ustalmy $\varepsilon>0$, zgodną metrykę zupełną $d$
i przeliczalny zbiór gęsty $\{x_j:j\ge1\}$.
Dla każdego $m\ge1$ oznaczmy
$V_{m,k}=\bigcup_{j=1}^k B(x_j,2^{-m})$.
Dla ustalonego $m$ zbiory $V_{m,k}$ rosną do $X$.
Wykażemy, że istnieje $k_m$ takie, że
\begin{equation}\label{eq:prokhorov-uniform-balls-22}
\inf_{\rho\in\mathcal F}\rho(V_{m,k_m})
>1-\varepsilon2^{-m-1}.
\end{equation}
Gdyby tak nie było, dla każdego $k$ wybralibyśmy
$\rho_k\in\mathcal F$ z
$\rho_k(V_{m,k})\le1-\varepsilon2^{-m-1}$.
Z założenia (ii) ciąg ten ma podciąg
$\rho_{k_l}\Rightarrow\rho$ z granicą probabilistyczną.
Ponieważ $V_{m,k}\uparrow X$, istnieje $k_0$ takie,
że $\rho(V_{m,k_0})>1-\varepsilon2^{-m-1}$.
Portmanteau dla otwartego $V_{m,k_0}$ daje
$\rho_{k_l}(V_{m,k_0})>1-\varepsilon2^{-m-1}$
dla dostatecznie dużych $l$. Ponieważ $k_l\ge k_0$,
monotoniczność wymusza wówczas
$\rho_{k_l}(V_{m,k_l})>1-\varepsilon2^{-m-1}$,
sprzecznie z konstrukcją $\rho_{k_l}$.
Dowiedliśmy \eqref{eq:prokhorov-uniform-balls-22}
bez odwoływania się do metryzowalności
przestrzeni prawdopodobieństw.

Niech $F_m=\bigcup_{j\le k_m}\overline B(x_j,2^{-m})$
i $K=\bigcap_{m\ge1}F_m$. Zbiór $K$ jest domknięty,
a dla każdego $m$ mieści się w skończonej sumie kul
o promieniu $2^{-m}$, więc jest całkowicie ograniczony.
Zupełność $X$ zapewnia zwartość $K$.
Z \eqref{eq:prokhorov-uniform-balls-22} dla wszystkich
$\rho\in\mathcal F$ uzyskujemy
\[
 \rho(X\setminus K)\le\sum_{m\ge1}\rho(X\setminus F_m)
 <\sum_{m\ge1}\varepsilon2^{-m-1}<\varepsilon.
\]
Jest to jednostajna ciasność $\mathcal F$.
\end{proof}

\begin{uwaga}[Metryzowalność a dowód Prochorowa]
Na przestrzeni polskiej topologia zbieżności wąskiej
na $\mathcal P(X)$ jest metryzowalna, na przykład przez
odległość opartą na ograniczonych funkcjach Lipschitzowskich.
Nie było jednak potrzeby przyjmowania tego wyniku
w dowodzie twierdzenia~\ref{thm:prokhorov-22}:
własność wyrażona przez podciągi i argument ze
wzrastającymi skończonymi pokryciami kulowymi
dały bezpośrednią konstrukcję wspólnego zwartego zbioru.
\end{uwaga}

\begin{wniosek}[Ciasność z oszacowania momentów]
\label{cor:coercive-tight-22}
Niech $X$ będzie przestrzenią polską i
$V:X\to[0,+\infty]$ funkcją dolnie półciągłą,
której zbiory podpoziomicowe
$K_R=\{x\in X:V(x)\le R\}$ są zwarte dla każdego
$R<\infty$. Jeśli rodzina $\mathcal F\subseteq\mathcal P(X)$
spełnia
$\sup_{\rho\in\mathcal F}\int_X V(x)\,\mathrm d\rho(x)
\le C<\infty$, to jest jednostajnie ciasna.
W szczególności na $\mathbb R^d$ wystarcza jednostajne
ograniczenie dodatniego momentu rzędu $p>0$.
\end{wniosek}
\begin{proof}
Dla $R>0$ nierówność Czebyszewa daje
\[
 \rho(X\setminus K_R)=\rho(\{V>R\})
 \le R^{-1}\int_X V(x)\,\mathrm d\rho(x)\le C/R.
\]
Wybierając $R>C/\varepsilon$ otrzymujemy
$\inf_{\rho\in\mathcal F}\rho(K_R)>1-\varepsilon$.
Na $\mathbb R^d$ funkcja $V(x)=|x|^p$ jest ciągła,
a jej zbiory podpoziomicowe są zwarte.
\end{proof}

\section{Następstwa i kryteria praktyczne}
\begin{wniosek}[Zbieżność mglista bez ucieczki masy]
\label{cor:vague-narrow-22}
Niech $X$ będzie lokalnie zwartą ośrodkową
przestrzenią metryczną. Jeśli $\mu_n,\mu\in\mathcal P(X)$,
$\mu_n$ zbiega mgliście do $\mu$, to
$\mu_n\Rightarrow\mu$.
\end{wniosek}
\begin{proof}
Ustalmy $\varepsilon>0$. Z rozdziału
\ref{ch:radon-measures-20} miara $\mu$ jest ciasna.
Wybieramy zwarty $K$ o $\mu(K)>1-\varepsilon$
oraz $\chi\in C_c(X)$ z $0\le\chi\le1$,
$\chi|_K=1$. Wówczas
$\int_X\chi\,\mathrm d\mu>1-\varepsilon$,
a ze zbieżności mglistej
$\int_X\chi\,\mathrm d\mu_n>1-2\varepsilon$
dla dużych $n$. Dla $h\in C_b(X)$
funkcja $h\chi\in C_c(X)$, a więc
\[
\left|\int_Xh\,\mathrm d\mu_n-
\int_Xh\chi\,\mathrm d\mu_n\right|
\le\|h\|_\infty\int_X(1-\chi)\,\mathrm d\mu_n
<2\varepsilon\|h\|_\infty.
\]
Analogiczne oszacowanie zachodzi dla $\mu$,
z błędem $\varepsilon\|h\|_\infty$.
Zbieżność całek z $h\chi$ i dowolność $\varepsilon$
kończą dowód.
\end{proof}

\begin{przyklad}[Rozkłady jednostajnie skupione na zwartym przedziale]
Niech $X=\mathbb R$, a $\mu_n$ dowolnymi
prawdopodobieństwami spełniającymi
$\mu_n([-R,R])=1$ dla stałego $R<\infty$.
Rodzina $\{\mu_n\}$ jest jednostajnie ciasna.
Z twierdzenia Prochorowa każdy jej ciąg ma podciąg
wąsko zbieżny. Nie wynika stąd zbieżność całego ciągu:
na przykład miary $\delta_{-1},\delta_1,\delta_{-1},\ldots$
mają dwa różne punkty skupienia.
\end{przyklad}

\begin{przyklad}[Brak ciasności]
Rodzina $\{\delta_n:n\ge1\}\subseteq\mathcal P(\mathbb R)$
nie jest jednostajnie ciasna: dla każdego zwartego
$K\subseteq\mathbb R$ pewne $n$ nie należy do $K$.
Nie istnieje podciąg zbieżny wąsko do prawdopodobieństwa.
\end{przyklad}

\begin{uwaga}[Gdzie pojawi się ta teoria?]
Zbieżność wąska, ciasność i twierdzenie Prochorowa
stanowią narzędzia w teorii procesów stochastycznych,
rachunku prawdopodobieństwa, analizie wariacyjnej
oraz w zagadnieniach transportu optymalnego.
Dalszy rozdział skryptu poświęcimy mierze Haara
na grupach lokalnie zwartych, gdzie niezmienniczość
miary dostarczy kolejnego sposobu jej scharakteryzowania.
\end{uwaga}